\section{Belief Representation: Finite Derivations}
\label{sec:annex-p15-rationality-4-7}

\begin{annexitem}{annex:p15.r47.threshold}{thm:p15-r47-threshold}{Thresholds, filters and singleton-complement representation}
\begin{proof}
Write \(t=\min\{g(A),g(B)\}\). If \(t>0\), use threshold \(\beta=t\)
to obtain \(g(A\cap B)\geq t\); if \(t=0\), the same inequality follows
from nonnegativity. Monotonicity gives the reverse inequality. Conversely
minitivity preserves acceptance at every threshold. A single threshold
would give only its own level set's closure, not this equality for all levels.

For Boolean \(N\) on a finite space, its accepted family is finite,
contains \(S\), and is closed under intersection. Its intersection \(E\)
is therefore itself accepted and cannot be empty, since \(N(\varnothing)=0\).
Every accepted event contains \(E\), and every event containing \(E\)
is accepted by monotonicity. For real-valued \(N\), express
\(A=\bigcap_{s\notin A}(S\setminus\{s\})\). Repeated minitivity gives
the displayed formula, including \(A=S\) via the empty minimum one.
Taking complements gives the maximum formula for \(\Pi\). Applying
the minimum formula to the empty event gives \(\min_s\iota(s)=0\),
hence \(\max_s\pi(s)=1\). Finally minitivity at \(A\cap A^c\)
gives \(\min(N(A),N(A^c))=0\). Intersections of positive-necessity
events remain positive, and cannot include both an event and its complement.
\end{proof}\end{annexitem}

\begin{annexitem}{annex:p15.r47.big-stepped}{thm:p15-r47-big-stepped}{Dominant contextual states and the lottery contrast}
\begin{proof}
Let \(s_i\) be the earliest positive-mass member of \(C\). All remaining
members have later indices. The big-step inequality gives
\(p_i>\sum_{s_j\in C,j>i}p_j\), unless \(i=n\), when the remaining
sum is zero and \(p_i>0\). Thus its conditional mass is greater than
one half. An event containing it has conditional mass greater than one
half, while an event missing it has less than one half. Since conditional
probabilities of an event and its complement sum to one, this is precisely
the displayed comparative acceptance test. Two accepted events contain
the same dominant state, so their intersection is accepted. Division
requires \(P(C)>0\), exactly the stated condition.

For the uniform lottery, choose \(n>1/(1-\beta)\) for any fixed
\(\beta<1\). Each losing event has probability \(1-1/n>\beta\),
but no state belongs to their joint intersection. This verifies the source
contrast without claiming that a threshold of one has the same failure.
\end{proof}\end{annexitem}

\begin{annexitem}{annex:p15.r47.ranking}{thm:p15-r47-ranking}{The static transform with the correct complement orientation}
\begin{proof}
Since exponentiation by \(2^{-t}\) reverses order, the minimum rank on
an event becomes the maximum singleton possibility. The zero minimum
rank gives maximum possibility one, and infinite empty rank gives empty
possibility zero. Event union takes the minimum of the two ranks, hence
the maximum of their possibilities. Complement duality then gives
necessity normalization and minitivity. The existing finite nested-focal
proof in \cref{thm:p15-ip-possibility-bridge} applies to this very normalized
distribution and gives its coherent lower-envelope representation.
No conditional, betting, or independence assertion occurs in that static
construction. In particular the printed alternative
\(N_\kappa(A^c)=2^{-\kappa(A)}\) would give \(N_\kappa(\varnothing)=1\)
at \(A=S\); it cannot be the normalized necessity used in the chapter.
The corrected complement formula is the one obtained from the actual
possibility/necessity duality proved above.
\end{proof}\end{annexitem}

\begin{annexitem}{annex:p15.r47.finite-modal}{thm:p15-r47-finite-modal}{Exhaustive epistemic-state evaluation and strict comparison}
\begin{proof}
There are \(r^n\) maps from \(S\) to the dual finite scale. Exactly
\((r-1)^n\) omit the value one, so excluding them leaves the asserted
number of normalized models. In each model compute the finite maximum
on every objective formula's truth set and its complement. This fixes
each graded atom; recursively apply the classical outer Boolean truth
functions to every premise and the conclusion. A counterexample is a
listed model satisfying all premises and falsifying the conclusion.
Every possible model is listed, so existence or absence of such a model
is exactly the stated semantic consequence, and the finite search terminates.

If \(N(p)>N(q)\), its positive value \(\lambda=N(p)\) is a grade in
\(L\); the corresponding disjunct is true. Conversely any true disjunct
gives \(N(p)\geq\lambda>N(q)\). Finiteness also gives
\(N(p)<\lambda\) exactly when \(N(p)\leq s(\lambda)\), explaining
the predecessor shift. None of these truth calculations establishes an
axiomatic completeness theorem for an unspecified inference system.
\end{proof}\end{annexitem}

\begin{annexitem}{annex:p15.r47.ignorance}{thm:p15-r47-ignorance}{Refinement and logarithmic reparametrization}
\begin{proof}
The first partition gives one of two equal probabilities to life. The
refined partition gives two of three, so \(2/3-1/2=1/6\). For the
continuous example use \(x=e^y\) and \(dx/dy=e^y\). Multiplying the
constant density \(1/(b-a)\) by that derivative gives the stated
nonconstant density. Conversely \(dy/dx=1/x\) multiplies the constant
log-scale density \(1/(\log b-\log a)\), giving
\(1/(x\log(b/a))\). The positive interval ensures both transforms are
defined and their densities integrate to one. These incompatible
uniformities cannot represent one refinement- and coordinate-invariant law.
\end{proof}\end{annexitem}

\begin{annexitem}{annex:p15.r47.focal}{thm:p15-r47-focal}{Nested versus overlapping focal sets and complete monotonicity}
\begin{proof}
Order the positive focal sets in an inclusion chain. The focal sets
contained in an event form an initial segment of that chain. The segment
contained in an intersection is the intersection of the two initial
segments, whose mass total is their smaller total. This proves minitivity.
Conversely if positive focal sets \(E,F\) are incomparable, their
intersection misses the positive mass of \(E\) counted in
\(\operatorname{Bel}(E)\), and misses the positive mass of \(F\)
counted in \(\operatorname{Bel}(F)\). Both beliefs therefore strictly
exceed the belief of their intersection, contradicting minitivity.

If disjoint positive focal sets exist, choose \(A=E\); both
\(\operatorname{Bel}(A)\) and \(\operatorname{Bel}(A^c)\) are positive.
Conversely, positive belief in both events supplies a positive focal
subset of each, hence two disjoint focal sets. Pairwise overlap is thus
exactly the weaker property. The three focal sets
\(\{1,2\},\{2,3\},\{1,3\}\), each of mass \(1/3\), are pairwise
overlapping with no common point, proving the stated limitation.

For a fixed focal set \(E\), apply finite inclusion--exclusion to the
Boolean predicates \(E\subseteq A_i\). Its alternating sum equals
one if \(E\) is contained in at least one \(A_i\), and zero otherwise.
This is no greater than the indicator \(E\subseteq\bigcup_iA_i\).
Multiply by \(m(E)\geq0\) and sum over all focal sets to get the full
displayed inequality for any finite \(k\). Singleton focal mass reduces
containment to membership and yields ordinary additivity. A mass-one
focal set gives its Boolean containment indicator. A chain gives the
minitive construction already proved. Incomparable positive focals show
why general belief functions cannot retain all-threshold adjunction.
\end{proof}\end{annexitem}

\begin{annexitem}{annex:p15.r47.combination}{thm:p15-r47-combination}{Nonempty-intersection normalization and simple supports}
\begin{proof}
The products \(m_1(E)m_2(F)\) are nonnegative and total one. The
nonempty-intersection products total \(K\); dividing their event-grouped
sums by positive \(K\) gives nonnegative masses totaling one, with empty
mass zero. For supports on the same \(E\), all four possible intersections
are nonempty. Only \((S,S)\) gives \(S\), with mass
\((1-p_1)(1-p_2)\); every other pair gives \(E\). Its mass is therefore
\(1-(1-p_1)(1-p_2)=p_1+p_2-p_1p_2\). Two disjoint mass-one sets have
no nonempty intersection product, so \(K=0\) and division is undefined.
The product calculation used the joint law of independent set-valued
sources. A dependent copy of one testimony has a different joint law:
for a perfectly copied simple support the set is \(E\) with probability
\(p\) and \(S\) with probability \(1-p\), so self-intersection leaves
reliability \(p\), not \(2p-p^2\). This is an explicit dependence failure.
\end{proof}\end{annexitem}

\begin{annexitem}{annex:p15.r47.pignistic}{thm:p15-r47-pignistic}{Unanimity-game allocation and the telescoping inverse}
\begin{proof}
For each nonempty focal set \(E\), distributing its mass equally over
its members gives total contribution \(m(E)\), so the displayed
nonnegative pignistic masses total one. The game
\(\operatorname{Bel}(A)\) is the weighted sum of unanimity games
\(u_E(A)=1\) iff \(E\subseteq A\). In a uniformly ordered permutation
of the players, a player contributes the marginal value one to \(u_E\)
only when it is the last member of \(E\) to enter the coalition.
Every member is equally likely to be last, with probability \(1/|E|\);
outsiders never contribute. The Shapley marginal average is linear in
the game values, giving exactly the pignistic formula.

For the inverse, \(p_i-p_{i+1}\geq0\). Summing the masses telescopes:
\(\sum_{i=1}^n i(p_i-p_{i+1})=\sum_{i=1}^np_i=1\).
The pignistic mass of \(s_j\) is
\(\sum_{i\geq j}i(p_i-p_{i+1})/i=p_j\).
The focal sets are nested, so the preceding theorem identifies its
belief with a necessity. Projection is not injective: on two states,
mass one on the whole space and equal singleton masses both project
to the uniform probability while having different singleton belief.
Thus the constructive inverse does not reconstruct a unique original
mass or prove an undefined least-information ordering.
\end{proof}\end{annexitem}

\begin{annexitem}{annex:p15.r47.price-boundaries}{thm:p15-r47-price-boundaries}{Sure loss and a strict maxmin-versus-dominance separation}
\begin{proof}
Buying one unit of each exclusive exhaustive indicator costs \(8/5\)
and returns exactly one in every state, yielding loss \(3/5\) surely.
The two isolated buy/sell gaps do not prevent it. For a credal envelope,
\(LP(f)=LP((f-g)+g)\geq LP(f-g)+LP(g)\), so a positive difference
lower expectation implies the asserted strict ordering of lower values.
For failure of the converse use the credal set of both point masses on
two states, \(f=(1,3)\), \(g=(0,4)\). Then
\(LP(f)=1>0=LP(g)\), but \(LP(f-g)=\min(1,-1)=-1\).
This is a precise stipulated decision model, not a report of Ellsberg
participants' preferences.
\end{proof}\end{annexitem}

\begin{annexitem}{annex:p15.r47.mobius}{thm:p15-r47-mobius}{Subset inversion and nonnegative focal masses}
\begin{proof}
Substituting the displayed inverse into \(\sum_{E\subseteq A}m(E)\)
gives, for each fixed \(B\subseteq A\), coefficient
\(\sum_{B\subseteq E\subseteq A}(-1)^{|E\setminus B|}\).
It is \((1-1)^{|A\setminus B|}\), equal to one for \(B=A\) and
zero otherwise. Thus the sum is \(b(A)\). Reversing the finite sums
likewise shows the inverse of a mass-generated function is that mass,
so it is unique. Normalization gives \(m(\varnothing)=b(\varnothing)=0\)
and \(\sum_E m(E)=b(S)=1\). Nonnegative inverse masses therefore
form a basic assignment, and every basic assignment necessarily has
that nonnegative inverse. No appeal to a general credal representation
can supply this missing nonnegativity condition.
\end{proof}\end{annexitem}

\begin{annexitem}{annex:p15.r47.likelihood}{thm:p15-r47-likelihood}{Finite singleton extension and normalization}
\begin{proof}
Every finite event is the union of its singletons. Repeated maxitivity
therefore forces its value to be the maximum singleton value; the empty
event has value zero. The maximum construction conversely has those
properties. If the largest likelihood \(M\) is positive, the rescaled
singleton scores lie in \([0,1]\) and at least one equals one, proving
normalization. If \(M=0\), all scores are zero and no ratio by \(M\)
is defined. A common positive multiplier changes neither their ordering
nor the normalized ratios. This establishes only the explicitly required
maxitive extension, not a posterior probability without prior information.
\end{proof}\end{annexitem}
